\section{Conclusão}
\label{sec:conclusao}

Este artigo investigou quem está exposto à inteligência artificial no mercado de trabalho brasileiro e desenvolveu um arcabouço teórico para fundamentar os padrões empíricos observados. A exposição à tecnologia --- mensurada pelo índice de \citet{eloundou2024gpts} aplicado aos microdados da PNAD Contínua --- concentra-se no núcleo formal e de maior capital humano: após controlar para um vetor completo de características mincerianas (idade, idade ao quadrado, gênero e raça), cada ano adicional de escolaridade associa-se a 0{,}23 ponto a mais de exposição tecnológica em regressões no nível individual com erros-padrão agrupados por ocupação ($N = 227.629$). Por outro lado, as ocupações que concentram a maioria dos trabalhadores brasileiros (serviços domésticos, construção estrutural e agropecuária) permanecem com exposição praticamente nula.

A relação entre exposição e formalidade é substancialmente mediada pela escolaridade, e o gradiente é mais íngreme nas Unidades da Federação de setor formal mais profundo. O modelo teórico de designação de tarefas com adoção corporativa versus individual de IA explica esse conjunto de fatos: a informalidade atua como uma barreira protetiva não por vedar o acesso a ferramentas pontuais, mas por limitar a capacidade de extração de ganhos organizacionais de produtividade que justificam o custo de formalização.

Para a formulação de políticas públicas, os resultados redefinem as prioridades de adaptação tecnológica. Se a informalidade amortece a exposição, as agendas de formalização do trabalho e de capacitação digital interagem de modo estratégico: formalizar sem capacitação pode expor trabalhadores a choques tecnológicos desprovidos de ferramentas adaptativas, ao passo que capacitar isoladamente sem expansão do setor formal restringe os ganhos de escala corporativa. O público prioritário para programas de requalificação e adaptação à IA não reside na base vulnerável do trabalho informal, mas no segmento médio-formalizado (técnicos, analistas e trabalhadores de apoio administrativo), hoje o mais exposto às transformações da automação informacional.
